% ============================================================
%  chapitres/introduction.tex  —  Chapitre 1 : Introduction
% ============================================================

\chapter{Introduction}

% ------------------------------------------------------------
\section{Présentation de l'organisation d'accueil}

Dona's Tech Estate est une agence de solutions numériques fondée sur la conviction que
l'ambition, et non le temps ou les antécédents de crédit, doit dicter le potentiel économique
de chaque entrepreneur. Positionnée comme le \textit{«~partenaire technique~»} de
l'\textit{«~entrepreneur parallèle~»} moderne, l'organisation s'adresse en priorité aux
immigrants instruits et aux professionnels ambitieux souhaitant lancer une entreprise durable
sans sacrifier leur emploi principal.

L'agence se distingue par une approche globale et clé en main, résumée sous le concept
d'\textit{«~entreprise en boîte~»} : plutôt que de fournir des services isolés (site web,
marketing), elle déploie un écosystème opérationnel complet couvrant la conformité légale,
l'automatisation des processus, le financement et la transformation numérique.

% ------------------------------------------------------------
\subsection{Activités et domaines d'expertise}

Les activités de Dona's Tech Estate s'articulent autour de quatre axes principaux :

\begin{itemize}
  \item \textbf{Blindage administratif :} gestion de la conformité légale complexe
        (Loi~25, Loi~96, incorporation fédérale et provinciale) afin que les clients
        puissent opérer en toute conformité et sans risque juridique ;

  \item \textbf{Automatisation et CRM (\textit{Command Center}) :} mise en place
        d'espaces de travail numériques automatisés en marque blanche~— hébergement web,
        système de réservation automatisé, CRM intégré~— jouant le rôle d'un personnel
        virtuel ;

  \item \textbf{Navigation pour le capital et le financement :} élaboration de feuilles
        de route pour l'obtention de subventions et de microcrédits canadiens, grâce à
        des modèles de plans d'affaires éprouvés contournant les barrières traditionnelles
        du crédit ;

  \item \textbf{Transformation numérique :} développement web axé sur la performance,
        référencement naturel (SEO) et intégration de flux de travail assistés par
        l'intelligence artificielle, afin d'assurer une compétitivité de niveau entreprise
        dès le lancement.
\end{itemize}

% ------------------------------------------------------------
\subsection{Produits et solutions}

L'organisation développe et maintient deux applications principales :

\begin{itemize}
  \item \textbf{SmartClinic :} plateforme web de gestion clinique permettant la prise
        de rendez-vous, la gestion des dossiers patients et la facturation ;
  \item \textbf{Plateforme OBNL :} outil de gestion destiné aux organismes à but non
        lucratif, couvrant la gestion des bénévoles, des donateurs et des activités.
\end{itemize}

% ------------------------------------------------------------
\subsection{Présence et clientèle}

Dona's Tech Estate opère principalement sur le marché canadien, avec un ancrage fort dans
la région de Gatineau\,/\,Ottawa. Sa clientèle est composée de startups en phase de
lancement ainsi que d'entreprises établies souhaitant moderniser leur infrastructure
numérique. L'agence intervient aussi bien auprès de structures à vocation commerciale que
d'organismes à but non lucratif (OBNL), ce qui explique la double nature de ses applications
phares.

% ------------------------------------------------------------
\subsection{Mission et vision}

La mission de Dona's Tech Estate est de combler le fossé entre l'emploi professionnel
salarié et la propriété d'une entreprise, en fournissant l'infrastructure automatisée,
conforme et professionnelle nécessaire à un lancement serein. Comme l'exprime sa PDG,
Dona~H :

\begin{quote}
  \itshape
  «~Nous aimons ce que nous faisons et nous aimons aider les autres à réussir dans ce
  qu'ils aiment faire.~»
\end{quote}

Sa vision repose sur la conviction que le \textit{«~side hustle~»} ne doit pas être un
fardeau chaotique, mais un actif stratégique et automatisé. L'organisation s'est donnée
pour mission de contribuer à l'indépendance financière de la main-d'œuvre immigrée, en
transformant le revenu actif en un patrimoine durable et générationnel.

% ------------------------------------------------------------
\subsection{Contexte du stage}

Ce stage s'inscrit dans le cadre de ma formation académique et a été réalisé au sein du
département de développement logiciel de Dona's Tech Estate. Dans un contexte où les
cybermenaces ciblant les applications de santé et les organisations civiles se multiplient,
la direction a souhaité renforcer la posture de sécurité de ses deux applications phares
avant leur prochaine mise en production.

% ============================================================
\section{Structure du rapport}

Dans le chapitre~2, nous présentons l'état de l'art relatif aux trois domaines couverts
par le mandat : la modélisation des menaces et les politiques de contrôle d'accès,
l'intégration DevSecOps du codage sécurisé, et les tests de pénétration éthiques.

Dans le chapitre~3, nous décrivons le contexte et les objectifs précis du projet :
présentation des deux applications, identification des enjeux de sécurité et définition
du périmètre d'intervention.

Dans le chapitre~4, nous exposons la démarche de modélisation des menaces appliquée à
SmartClinic et à la Plateforme OBNL à l'aide de la méthodologie STRIDE et des diagrammes
de flux de données (DFD).

Dans le chapitre~5, nous définissons les politiques de contrôle d'accès retenues pour les
deux applications en justifiant le choix du modèle RBAC et les mécanismes d'authentification
renforcée.

Dans le chapitre~6, nous décrivons la mise en place des processus DevSecOps : intégration
des outils SAST, DAST et SCA dans le pipeline CI/CD, règles de codage sécurisé et gestion
des dépendances.

Dans le chapitre~7, nous présentons les tests de pénétration éthiques réalisés sur les
environnements de test, la méthodologie suivie, les vulnérabilités identifiées et leur
niveau de criticité (CVSS).

Dans le chapitre~8, nous analysons les résultats globaux et discutons l'amélioration de la
posture de sécurité des deux applications.

Dans le chapitre~9, nous formulons des recommandations concrètes à destination des équipes
de développement et de la direction.
